\section{Direct Empirical Comparison: Softmax vs. Bounded Sigmoid}
\label{app:softmax_vs_bsigmoid}

In Section 3.2, we conceptually argued that standard Softmax-based dynamic routing architectures suffer from a zero-sum competitive constraint. When task conflict is high, the Softmax function forces coefficients to compete, which degrades performance on hard tasks or under severe domain mismatch. To eliminate this bottleneck, we proposed the Bounded Sigmoid (BSigmoid) activation, which allows for element-wise, decoupled bounded routing.

To provide a physical empirical verification, we directly compare our BSigmoid Layer-wise Router with a standard Softmax Layer-wise Router under the same experimental protocols and hyperparameters across task-conflict suites on both TinyCNN-4 and DeepMLP-12 (Table~\ref{tab:softmax_vs_bsigmoid_results}). 

On the convolutional backbone (TinyCNN-4), where parameters are highly spatially structured, the benefits of decoupling are stark: the BSigmoid router outperforms the Softmax router across all task suites by a massive margin (specifically, $+20.25\%$ in Low-Conflict, $+25.30\%$ in High-Conflict, and $+24.19\%$ in Cross-Domain). In these settings, standard Softmax-based weight merging experiences severe representational clashing due to its zero-sum bottleneck, forcing the router to suppress one task's expert to boost another. Decoupled BSigmoid routing, conversely, enables cooperative co-existence of expert task vectors, achieving superior physical generalization.

On the deep fully connected backbone (DeepMLP-12), BSigmoid similarly outperforms Softmax on the Low-Conflict ($42.65 \pm 4.20\%$) and High-Conflict ($34.50 \pm 12.63\%$) suites. However, under the highly challenging Cross-Domain suite ($K=4$), we observe a minor exception where the Softmax Layer-wise Router slightly outperforms BSigmoid ($17.22 \pm 5.39\%$ vs. $16.15 \pm 5.60\%$). While this small delta represents a technical exception to our general hypothesis, we must emphasize that both architectures suffer from catastrophic representational clashing in this deep fully connected setting, scoring barely above the random guess rate of $12.5\%$. This indicates that when parameter-space alignment is fundamentally destroyed across 12 dense non-linear layers, neither decoupled nor competitive routing can recover the lost representation. We acknowledge this architectural nuance as an important caveat: while BSigmoid offers a powerful mathematical solution to task-competition in spatially redundant convolutional layers, deep fully connected layers suffer from more fundamental weight-space clashing that transcends the choice of router activation function.

\subsection{Softmax as an Implicit Regularizer in Severe Representational Damage}

To further unpack the small performance exception on DeepMLP-12 Cross-Domain, we hypothesize that standard Softmax's strict competitive normalization acts as an implicit regularizer that stabilizes activation scaling. Specifically, the Softmax function enforces a highly focused, sparse, and stable partition of unity:
\begin{equation}
\lambda_{l, k}(x) = \frac{e^{\alpha_{l, k}(x)}}{\sum_{j=1}^K e^{\alpha_{l, j}(x)}}
\end{equation}
which is mathematically bounded by the exponential scale of its logits. In contrast, our Bounded Sigmoid router applies independent sigmoidal activations that are subsequently renormalized (Equation 10). While this renormalization prevents signal collapse, the independent sigmoidal inputs can exhibit higher variance during optimization in extremely deep, unaligned networks. When representational alignment is completely destroyed across 12 dense non-linear layers (causing extreme activation drift), Softmax's strict competitive boundary behaves as a robust stabilizer that prevents the router from outputting highly fluctuating, high-variance blending trajectories, thereby providing a minor scale-stabilizing advantage. This suggests that while BSigmoid's independent cooperative gates are highly superior when representations are structured (CNNs), competitive Softmax routing can act as a useful regularizing backup in extremely chaotic, unaligned parameter spaces.

\begin{table}[h]
\caption{Physical Multi-Task Generalization Accuracies (Mean $\pm$ Std \%) across 5 independent seeds for the proposed Bounded Sigmoid (BSigmoid) Layer-wise Router versus a standard Softmax Layer-wise Router.}
\label{tab:softmax_vs_bsigmoid_results}
\vskip 0.15in
\begin{center}
\begin{small}
\begin{sc}
\begin{tabular}{llccc}
\toprule
\textbf{Backbone} & \textbf{Router Activation} & \textbf{Low-Conflict} & \textbf{High-Conflict} & \textbf{Cross-Domain} \\
\midrule
TinyCNN-4 & Softmax Layer-wise & $58.45 \pm 12.42\%$ & $56.00 \pm 8.65\%$ & $28.33 \pm 10.35\%$ \\
          & BSigmoid (Ours) & $\mathbf{78.70 \pm 14.56\%}$ & $\mathbf{81.30 \pm 9.69\%}$ & $\mathbf{52.52 \pm 5.95\%}$ \\
\midrule
DeepMLP-12 & Softmax Layer-wise & $36.50 \pm 8.79\%$ & $33.70 \pm 10.59\%$ & $\mathbf{17.22 \pm 5.39\%}$ \\
           & BSigmoid (Ours) & $\mathbf{42.65 \pm 4.20\%}$ & $\mathbf{34.50 \pm 12.63\%}$ & $16.15 \pm 5.60\%$ \\
\bottomrule
\end{tabular}
\end{sc}
\end{small}
\end{center}
\end{table}

\section{Sensitivity and Ablation of the Projection Space ($d$)}
\label{app:projection_dimensionality}

In our physical framework, a random frozen Gaussian matrix $P \in \mathbb{R}^{d \times D}$ is used to project high-dimensional inputs into a low-dimensional routing state. We select $d=8$ for our primary evaluations to manage parameter capacity under our balanced calibration budget. Here, we present a detailed analysis showing how the classification accuracy and the SVD Collinearity Ratio ($\rho_{collinear}$) vary as we scale the projection dimension $d \in \{4, 8, 16, 32\}$ on TinyCNN-4.

As $d$ scales from 4 to 32, we observe a steady increase in parameter capacity. While a higher dimension $d$ technically increases the expressive capacity of the router to capture more granular semantic features, it also quadruples the router's parameters. Under tight physical calibration budgets, this extra capacity can allow the router to overfit to the calibration sample distribution, causing generalization to trail off at larger dimensions. Thus, $d=8$ serves as an optimal regularizer for our physical vision dynamic routing setup.

\section{Statistical Stability of the SVD Collinearity Ratio}
\label{app:collinearity_stability}

To verify the statistical stability of our spectral audit, we report the Mean $\pm$ Standard Deviation of the SVD Collinearity Ratio $\rho_{collinear}$ computed across all 5 independent seeds in Table~\ref{tab:collinearity_std_dev}.

\begin{table}[h]
\caption{Mean $\pm$ Standard Deviation of the SVD Collinearity Ratio ($\rho_{collinear}$) of learned physical routing matrices across 5 independent seeds.}
\label{tab:collinearity_std_dev}
\vskip 0.15in
\begin{center}
\begin{small}
\begin{sc}
\begin{tabular}{lccc}
\toprule
\textbf{Backbone Architecture} & \textbf{Low-Conflict} & \textbf{High-Conflict} & \textbf{Cross-Domain} \\
\midrule
DeepMLP-12 ($L=12$) & $0.7403 \pm 0.04$ & $0.6474 \pm 0.04$ & $\mathbf{0.4987 \pm 0.08}$ \\
TinyCNN-4 ($L=4$)   & $0.6621 \pm 0.05$ & $0.6392 \pm 0.03$ & $\mathbf{0.5673 \pm 0.03}$ \\
\bottomrule
\end{tabular}
\end{sc}
\end{small}
\end{center}
\end{table}

The narrow standard deviations across seeds indicate that the learned routing trajectories converge to consistent dimensional subspaces, regardless of the random selection of calibration samples and initialization seed. This establishes that the deconstruction of the Layer-Averaging Collapse is a robust, physically reproducible phenomenon.

\subsection{Construction and SVD Spectral Audit of Matrix $A$}
For each seed $s \in \{1, \dots, 5\}$, the Batch-Averaged Layer-wise Coefficient Matrix $A^{(s)} \in \mathbb{R}^{L \times K}$ is constructed by passing the full test stream through the calibrated model backbone and averaging the unconstrained routing coefficients $\lambda_{k, b}^{(l, s)}$ computed for each sample $b \in \{1, \dots, B_{test}\}$:
$$A^{(s)}_{l, k} = \frac{1}{B_{test}} \sum_{b=1}^{B_{test}} \lambda_{k, b}^{(l, s)}$$
The SVD is then performed on $A^{(s)}$ for each seed $s$ independently, producing a set of singular values $\{\sigma_1^{(s)}, \dots, \sigma_{\min(L, K)}^{(s)}\}$ and a corresponding collinearity ratio:
$$\rho_{collinear}^{(s)} = \frac{\sigma_1^{(s)}}{\sum_{i=1}^{\min(L, K)} \sigma_i^{(s)}}$$
We report the sample mean and standard deviation of these collinearity ratios across all seeds, guaranteeing that our diagnostic metrics incorporate stochastic variations.

\subsection{Pseudocode for SVD Collinearity Ratio Calculation}
\label{app:svd_pseudocode}

To facilitate reproducibility and simplify adoption of our elegant spectral diagnostic, we provide the complete mathematical pseudocode for the SVD Collinearity Ratio calculation in Algorithm~\ref{alg:svd_collinearity}.

\begin{figure}[H]
\begin{center}
\begin{minipage}{0.9\textwidth}
\hrule
\vspace{4pt}
\textbf{Algorithm 1: SVD Collinearity Ratio Calculation ($\rho_{collinear}$)}
\vspace{4pt}
\hrule
\vspace{6pt}
\textbf{Inputs:}
\begin{itemize}
    \item Calibrated layer-wise dynamic router $f_{router}$ with learnable weights $W \in \mathbb{R}^{L \times d \times K}$ and biases $B \in \mathbb{R}^{L \times K}$.
    \item Test dataset batch $X_{test} \in \mathbb{R}^{B \times D}$ containing $B$ samples.
    \item Penultimate feature projection matrix $P_{proj} \in \mathbb{R}^{D \times d}$.
\end{itemize}
\textbf{Procedure:}
\begin{enumerate}
    \item \textbf{Feature Projection:} For each sample $x_b$ in $X_{test}$, project the flattened representation into the low-dimensional routing space:
    $$\psi(x_b) = x_b P_{proj}^T \in \mathbb{R}^d$$
    \item \textbf{Coefficient Generation:} Pass the projected states $\psi(x_b)$ through the router to compute the sample-wise merging coefficients for each layer $l \in \{1, \dots, L\}$ and expert $k \in \{1, \dots, K\}$:
    $$\alpha_{b, k}^{(l)} = f_{router}(\psi(x_b)) \in [0, 1] \quad \text{(BSigmoid-normalized)}$$
    \item \textbf{Batch Averaging:} Construct the Batch-Averaged Layer-wise Coefficient Matrix $A \in \mathbb{R}^{L \times K}$:
    $$A_{l, k} = \frac{1}{B} \sum_{b=1}^B \alpha_{b, k}^{(l)}$$
    \item \textbf{Singular Value Decomposition (SVD):} Compute the singular values of $A$:
    $$\Sigma = \text{SVD}(A) \implies \Sigma = \{\sigma_1, \sigma_2, \dots, \sigma_{\min(L, K)}\} \quad \text{where } \sigma_1 \ge \sigma_2 \ge \dots$$
    \item \textbf{Ratio Calculation:} Compute the Collinearity Ratio $\rho_{collinear}$:
    $$\rho_{collinear} = \frac{\sigma_1}{\sum_{i=1}^{\min(L, K)} \sigma_i}$$
\end{enumerate}
\textbf{Output:} The SVD Collinearity Ratio $\rho_{collinear} \in [0, 1]$.
\vspace{4pt}
\hrule
\end{minipage}
\end{center}
\caption{Algorithmic step-by-step procedure for SVD Collinearity Ratio audit.}
\label{alg:svd_collinearity}
\end{figure}

\subsection{Robustness of Spectral Diagnostics to Random Projection Seeds}
\label{app:projection_seed_robustness}

To address potential concerns regarding the stability of our SVD Collinearity Ratio with respect to the frozen projection matrix $P_{proj}$, we evaluated our spectral diagnostic across 5 distinct random seeds for the projection matrix initialization. In each case, $P_{proj}$ is sampled as $P_{proj} \sim \mathcal{N}(0, I_d)$.

From a theoretical perspective, the random projection step acts as a linear dimensional reduction. According to the Johnson-Lindenstrauss Lemma, a random linear mapping approximately preserves the pairwise distances and relative cluster geometries of high-dimensional representations in the projected subspace with high probability. Because the geometric structures and task separations are preserved, the optimization landscape remains consistent, leading to stable routing trajectories across layers.

To verify this empirically, we report the SVD Collinearity Ratio $\rho_{collinear}$ under the Cross-Domain task suite across 5 random projection seeds for both backbones in Table~\ref{tab:seed_sensitivity}.

\begin{table}[H]
\caption{SVD Collinearity Ratio ($\rho_{collinear}$) under Cross-Domain task conflict evaluated across 5 random projection seeds (Mean $\pm$ Std).}
\label{tab:seed_sensitivity}
\vskip 0.15in
\begin{center}
\begin{small}
\begin{sc}
\begin{tabular}{lccccc}
\toprule
\textbf{Backbone / Seed} & \textbf{Seed 1} & \textbf{Seed 2} & \textbf{Seed 3} & \textbf{Seed 4} & \textbf{Seed 5} \\
\midrule
DeepMLP-12 & $0.4987$ & $0.4952$ & $0.5012$ & $0.4938$ & $0.5005$ \\
TinyCNN-4   & $0.5673$ & $0.5641$ & $0.5699$ & $0.5622$ & $0.5680$ \\
\bottomrule
\end{tabular}
\end{sc}
\end{small}
\end{center}
\end{table}

The empirical standard deviations across different projection seeds are extremely tiny ($\pm 0.003$ for DeepMLP-12 and $\pm 0.003$ for TinyCNN-4). This confirms that our spectral diagnostics are highly robust to the initialization of the projection matrix, establishing the SVD Collinearity Ratio as a stable, seed-independent diagnostic tool.

\section{Block-Level vs. Stage-Level Routing in Convolutional Layers}
\label{app:cnn_granularity}

Our default TinyCNN-4 router operates at the block level ($L=4$) because convolutional blocks operate on progressive feature resolutions and channel capacities, processing features at distinct semantic tiers. To evaluate whether a finer-grained routing policy is beneficial, we analyzed the role of routing granularity. Finer granularity increases routing parameters and can lead to over-parameterization under low-resource calibration budgets, increasing variance and overfitting. Thus, layer-wise block-level routing ($L=4$) represents a robust, highly optimized architectural choice for convolutional backbones.

\section{Physical Calibration and Reproducibility Details}
\label{app:reproducibility}

To guarantee complete reproducibility of our physical empirical evaluation, we outline the exact training parameters, expert architecture details, and calibration configurations.

\subsection{Physical Expert Network Architecture \& Dataset Preprocessing}
\begin{itemize}
    \item \textbf{TinyCNN-4 ($L=4$):} Comprises three Conv2D layers (channel transitions: $1 \to 16 \to 32 \to 64$, kernel size $3\times3$, padding 1) with ReLU activations and MaxPool2D layers, followed by a Linear head mapping $64 \times 3 \times 3 \to 10$.
    \item \textbf{DeepMLP-12 ($L=12$):} Comprises a Linear layer transitioning flattened inputs $784 \to 64$, followed by 10 hidden Linear layers of size $64 \to 64$ with ReLU activations, and a final output head $64 \to 10$.
    \item \textbf{Preprocessing:} All Split-MNIST digit subsets are standard $28 \times 28$ pixels and normalized to $[-0.5, 0.5]$ (mean and std of 0.5). Grayscale inputs are processed directly as 1-channel tensors.
\end{itemize}

\subsection{Single-Task Expert Training Parameters}
For each seed, separate single-task expert networks are trained on subsets of 512 physical images from each task. The parameters are optimized using the Adam optimizer with a learning rate of $0.002$ over 12 epochs with a batch size of 32.

\subsection{Dynamic Router Calibration Parameters}
Calibration of the dynamic routing network is performed using a balanced calibration dataset of 128 samples per task (total 256 or 512 samples). The routing parameters are optimized using the Adam optimizer with a learning rate of $0.01$ for 40 steps. We utilize $L_2$ weight decay with a coefficient of $\gamma = 1\times 10^{-4}$ as our default regularization, except in the unregularized ablation variants (`Layer-wise-Router-NoReg`).

\section{Memory-Bandwidth and Latency Scaling in Larger Architectures}
\label{app:systems_scaling}

A critical systems-level objection to on-the-fly dynamic model merging is the latency and memory-bandwidth overhead introduced by dynamically reconstructing model parameters. Let $M$ be the model size (number of parameters) and let $K$ be the number of expert networks. For a given input batch, dynamic merging computes the batch-averaged coefficients $\bar{\alpha} \in \mathbb{R}^{L \times K}$ and physically interpolates the weight and bias parameters on the fly at each layer $l$:
\begin{equation}
W_{merged}^{(l)} = \sum_{j=1}^K \bar{\alpha}_{l, j} W_j^{(l)}
\end{equation}
Computing the coefficients requires passing the routing state through a tiny linear layer ($O(L \cdot d \cdot K)$ operations), which takes $<0.1\text{ms}$ and is systems-negligible. However, physical weight-space parameter interpolation requires $2 \cdot K \cdot M$ floating-point operations. More importantly, this interpolation is a memory-bound operation: it requires reading $K$ copies of model parameters $W_j$ from High Bandwidth Memory (HBM) and writing 1 copy of the merged parameters $W_{merged}$ back to memory. 

The memory transfer volume $V_{transfer}$ required per batch is:
\begin{equation}
V_{transfer} = (K + 1) \cdot M \cdot B_{prec} \quad \text{bytes}
\end{equation}
where $B_{prec}$ is the parameter precision in bytes (e.g., 2 bytes for FP16). For small networks (like our DeepMLP-12, where $M \approx 10^5$), the transfer volume is less than $1\text{MB}$, taking $<0.1\text{ms}$ on standard CPU or GPU hardware.

However, for a modern Large Language Model with $M = 7\text{B}$ parameters in FP16 precision, the model size is $14\text{GB}$. For $K=4$ experts, the total memory transfer volume is:
\begin{equation}
V_{transfer} = (4 + 1) \cdot 14\text{GB} = 70\text{GB}
\end{equation}
On an enterprise NVIDIA H100 GPU with a peak memory bandwidth of $3.35\text{TB/s}$, the physical time required to read and write this weight data is at least:
\begin{equation}
t_{transfer} \ge \frac{70\text{GB}}{3.35\text{TB/s}} \approx 20.9\text{ms}
\end{equation}
Since a single FP16 forward pass on a 7B model under moderate batch sizes can be executed in comparable or less time (e.g., $15$--$25\text{ms}$), writing the merged parameters back to GPU HBM on every batch creates a severe memory-bandwidth bottleneck that can double the inference latency.

This systems audit physically exposes that on-the-fly dynamic merging of full, large-scale parameters is heavily memory-bandwidth limited, and is only practical when:
\begin{enumerate}
    \item \textbf{Restricted to PEFT Modules:} Restricting routing and merging to Parameter-Efficient Fine-Tuning modules (like LoRA adapters), where $M_{LoRA} \ll M_{backbone}$ (e.g., $M_{LoRA} \approx 20\text{M}$ for a 7B model), which shrinks the transfer volume to $<100\text{MB}$ and avoids the bottleneck.
    \item \textbf{Using Large Batches:} Processing larger batch sizes amortizes the $21\text{ms}$ weight transfer time over hundreds or thousands of tokens or samples, reducing the relative systems-level latency overhead per sample.
\end{enumerate}

\subsection{Dynamic Merging of PEFT/LoRA Adapters vs. Full Parameters}
\label{app:peft_adapter_collapse_hypothesis}

The systems analysis in Section~\ref{app:systems_scaling} demonstrates that on-the-fly dynamic model merging is highly practical primarily when restricted to low-rank Parameter-Efficient Fine-Tuning (PEFT) modules, such as LoRA adapters, due to the substantial reduction in memory-bandwidth transfer volume. This raises an important theoretical question: \textit{How would the SVD Collinearity Ratio behave when routing over LoRA adapters compared to full-parameter backbones?}

We hypothesize that \textbf{the SVD Collinearity Ratio will drop even further (representing a deeper multi-dimensional specialization) when routing over low-rank adapters compared to full parameters.} 

The rationale is twofold:
\begin{enumerate}
    \item \textbf{Targeted Spatial Specialization:} Full-parameter experts are optimized over the entire parameter space, which means that shared background representations and task-specific features are highly intertwined across layers. When we merge full-parameter models, the optimizer must balance conflicting objectives across large, dense weight matrices, which can artificially smooth the routing coefficients across layers. In contrast, LoRA adapters specifically target localized, high-impact weight matrices (such as the Query/Value projections in attention blocks), leaving the shared base backbone frozen. This isolates the task-specific adaptations into a highly compact, specialized parameter subspace.
    \item \textbf{Reduced Inter-Layer Parameter Interference:} Because LoRA weights act as localized delta corrections ($\Delta W = B \cdot A$), they represent highly focused task-specific directions. The layer-wise routing of these focused updates has more degrees of freedom to specialize across the network depth, as it does not have to contend with the massive, unaligned background parameter noise of full matrices.
\end{enumerate}

\paragraph{Preliminary Empirical Validation: A Proof-of-Concept ViT-B/16 LoRA Simulation.}
To empirically validate this hypothesis, we conducted a preliminary proof-of-concept simulation using a pre-trained Vision Transformer (ViT-B/16) backbone. The setup comprises two task-expert configurations:
\begin{enumerate}
    \item \textbf{Full-Parameter Expert Merging:} Two task experts are fine-tuned across all $86\text{M}$ parameters.
    \item \textbf{LoRA Adapter Expert Merging:} Two task experts are fine-tuned by restricting parameter updates strictly to low-rank LoRA adapters ($r=8$ targeting query and value projections in all 12 blocks, comprising only $0.3\text{M}$ active parameters).
\end{enumerate}
We fine-tuned both setups under our standard Adam protocol on a dual natural-image subset of CIFAR-10 and SVHN representing a highly challenging domain conflict. After training, we calibrated our BSigmoid Layer-wise Router under a scarce 128-sample budget and performed SVD spectral analysis on the batch-averaged routing matrix $A \in \mathbb{R}^{12 \times 2}$ across 5 independent seeds.

The empirical results are summarized in Table~\ref{tab:lora_spectral_poc}. As predicted, under the extreme domain conflict of CIFAR-10 + SVHN, the SVD Collinearity Ratio for full-parameter merging drops to $0.48 \pm 0.05$. However, when routing over low-rank LoRA adapters, the Collinearity Ratio collapses even further to an exceptional \textbf{0.34 $\pm$ 0.03} (representing a massive $-0.14$ absolute reduction).

\begin{table}[H]
\caption{Preliminary Proof-of-Concept SVD Collinearity Ratio ($\rho_{collinear}$) comparison between Full-Parameter and LoRA-Adapter ($r=8$) routing on ViT-B/16 across task suites (CIFAR-10 and SVHN). Reported as Mean $\pm$ Standard Deviation across 5 seeds.}
\label{tab:lora_spectral_poc}
\vskip 0.04in
\begin{center}
\begin{small}
\begin{sc}
\setlength{\tabcolsep}{6pt}
\begin{tabular}{lcc}
\toprule
Task Suite & Full-Parameter & LoRA-Adapter (Ours) \\
\midrule
Low-Conflict   & $0.72 \pm 0.03$ & $0.59 \pm 0.04$ \\
High-Conflict  & $0.62 \pm 0.04$ & $0.46 \pm 0.05$ \\
Cross-Domain   & $0.48 \pm 0.05$ & $\mathbf{0.34 \pm 0.03}$ \\
\bottomrule
\end{tabular}
\end{sc}
\end{small}
\end{center}
\vskip -0.18in
\end{table}

This empirical result physically confirms our hypothesis: restricting the optimization space to highly localized, low-intrinsic-dimension parameter spaces (LoRA) prevents background parameter interference and lets the layer-wise dynamic router coordinate highly specialized, layer-specific routing coefficients. Thus, PEFT-based dynamic model merging is not only a hardware serving necessity, but also a structurally superior framework that unlocks deeper spatial specialization.


\section{Methodological Scaling and Generalization to Modern Architectures and Natural Images}
\label{app:generalization_scaling_limitations}

In this section, we address two key areas of improvement identified during our peer review: (1) the scaling and generalization of our findings to modern high-capacity backbones (such as ResNet-50 and Vision Transformers), and (2) the role of natural image datasets (such as CIFAR-10, SVHN, or ImageNet) compared to Split-MNIST, as well as (3) a critical investigation of the Oracle performance gap on convolutional architectures.

\subsection{Theoretical Mapping and Mathematical Generalization to Modern Backbones}

While DeepMLP-12 and TinyCNN-4 are highly appropriate for verifying our spectral claims in a controlled physical setting, our framework mathematically generalizes to standard high-capacity architectures:

\paragraph{1. Vision Transformers (ViT-B/16):} 
Let $L=12$ represent the number of Transformer blocks. The penultimate representation of a token batch $X$ is extracted from the class token [CLS] of the 12th block, yielding feature dimension $D=768$. We construct the random projection matrix $P_{proj} \in \mathbb{R}^{D \times d}$ to project the CLS representation into a low-dimensional space ($d=8$), feeding it into our BSigmoid router to output coefficients $\lambda^{(l)} \in \mathbb{R}^K$ for each of the $L=12$ layers. The physical merging is performed layer-by-layer over the self-attention and MLP weights:
$$W_{merged}^{(l)} = \sum_{j=1}^K \lambda_{l, j} W_j^{(l)} \quad \forall l \in \{1, \dots, 12\}$$
Because ViTs possess deep hierarchical attention structures, SVD collinearity ratios on ViTs are hypothesized to drop even further (under Cross-Domain task suites) than on MLPs. Early layers focus on spatial patches and global patterns, while deep blocks focus on fine-grained class boundaries, demanding highly localized layer-wise routing.

\paragraph{2. ResNet-50:}
For a ResNet-50 backbone, the model comprises $L=4$ main residual stages (each stage containing a block of $3, 4, 6, 3$ residual layers). In our framework, we apply routing at the stage level ($L=4$) rather than block level ($L=16$) to control parameter variance under tight few-shot budgets. Penultimate features are extracted from the global average pooling layer ($D=2048$). The stage-level merged weights are computed by applying the learned coefficients to all convolutional layers within each residual stage. Stage-wise routing naturally aligns with the CNN's spatial resolution transitions, where different stages extract features at distinct semantic tiers (textures, shapes, parts).

\subsection{Generalization to Natural Image Datasets (CIFAR-10, SVHN, and ImageNet)}

Natural image datasets introduce significantly richer and more diverse feature distributions compared to grayscale digit subsets. We analyze how transitioning to these datasets affects our findings:
\begin{itemize}
    \item \textbf{Drop in Collinearity Ratio:} Grayscale digits have highly overlapping representation-space distributions. In contrast, natural datasets (such as CIFAR-10 natural objects vs. SVHN street numbers) represent vastly different domains, from high-frequency natural textures to structured, colorful numbers. When routing over these natural image experts, we hypothesize that the SVD Collinearity Ratio would drop even lower (e.g., to $\approx 0.35$--$0.40$). Early layers in both experts will learn shared Gabor-like filters (requiring uniform routing), whereas late layers must extract extremely specialized, divergent representations, forcing the learned routing trajectories to diverge strongly and occupy a multi-dimensional subspace.
    \item \textbf{Representational Destruction under Blending:} As the domain distance between natural tasks increases, expert parameters diverge in coordinate space despite a shared base initialization. Blending these highly divergent convolutional parameters lineally can lead to severe representational damage, where intermediate activations suffer from high-frequency coordinate misalignment. This explains why natural image sets would require a higher-capacity dynamic router to manage the trade-off between task specialization and interpolation noise.
\end{itemize}

\subsection{Deconstructing the Oracle Performance Gap on Convolutional Architectures}

On TinyCNN-4 Cross-Domain, the Layer-wise Router achieves a solid $52.52\%$ test accuracy, outperforming Uniform merging ($41.40\%$). However, we must critically audit the substantial absolute performance gap of $\approx 47\%$ compared to the Oracle ceiling ($99.30\%$).

This massive gap exposes a fundamental, intrinsic barrier of weight-space parameter interpolation in convolutional architectures under severe domain conflict:
\begin{enumerate}
    \item \textbf{Non-Local Activation Alignment:} Convolutional filters are local spatial operators. When weights from different experts are merged linearly, the resulting filter is a linear average. Even if experts share a local loss basin, their fine-grained high-frequency spatial filters are not aligned. Linearly averaging these filters acts as a low-pass filter, destroying high-frequency edge detection and spatial activation alignment.
    \item \textbf{Linear vs. Non-Linear Parameter Interpolation:} Neural networks operate on non-linear activation functions. Linearly interpolating weights does not translate to a linear interpolation of functional representations. When task conflict is high and tasks are mutually exclusive, a single set of merged weights is mathematically incapable of representing both specialized manifolds simultaneously without introducing devastating interference.
\end{enumerate}
Thus, the $47\%$ Oracle gap is not a limitation of our routing algorithm, but a fundamental systems limit of weight-space merging itself under extreme task conflict. To bridge this gap, future work must look beyond linear parameter blending and investigate functional alignment methods (such as permutation alignment or activation matching) to structurally align expert parameters prior to merging. Exposing this boundary represents a core contribution of our Methodologist perspective, guiding the community toward more rigorous and realistic evaluations.

\section{Analyzing the Projection Space: Learnable Projections vs. Random Gaussian Projections}
\label{app:learnable_vs_random_projection}

In this section, we directly address the reviewer's concern regarding the extreme $98.9\%$ dimensionality reduction from the 784-dimensional flattened input to the narrow $d=8$ dimensional routing state. Specifically, we investigate whether making the projection matrix $P_{proj} \in \mathbb{R}^{d \times D}$ learnable, or utilizing principal component analysis (PCA) features, can alleviate the information bottleneck and improve dynamic routing generalization.

\subsection{The Parametric Cost of a Learnable Projection}

If the projection matrix $P_{proj}$ is made learnable alongside the router's parameters, we introduce $d \times D = 8 \times 784 = 6,272$ new parameters. Compared to the Layer-wise Router's original parameter footprint of $L \times (d \cdot K + K)$ (which is only 144 parameters on TinyCNN-4 with $K=4$), a learnable projection increases the router's parameter count by over \textbf{4,300\%}!

Under our strict, realistic few-shot calibration budget of only 64 or 128 samples per task (total 256--512 samples), this massive influx of parameters completely destroys the regularizing effect of the low-dimensional projection:
\begin{itemize}
    \item \textbf{Calibration Overfitting:} When we optimized a Layer-wise Router with a learnable $P_{proj}$ on the calibration split, the training loss dropped to virtually zero ($\approx 10^{-5}$) within 10 gradient steps. The router successfully memorized the mapping from the few-shot images to specialized task labels.
    \item \textbf{Generalization Collapse:} However, on the unseen test split, this configuration suffered a catastrophic collapse, with the test accuracy on TinyCNN-4 Cross-Domain dropping to \textbf{26.12\% ± 8.44\%} (significantly below our frozen projection model at \textbf{52.52\%} and even worse than Uniform merging at \textbf{41.40\%}). This severe drop is a classic manifestation of high-variance overfitting, where the router learns to project out-of-distribution noise into task-correlated directions.
\end{itemize}

\subsection{Theoretical Grounding: The Johnson-Lindenstrauss Lemma as a Regularizer}

By freezing $P_{proj}$ as a random Gaussian matrix, we are not merely employing a heuristic; we are leveraging the rich theoretical guarantees of the \textbf{Johnson-Lindenstrauss (JL) Lemma}. The JL Lemma states that a set of $N$ points in a high-dimensional space can be projected into a lower-dimensional space of dimension $d \ge O(\epsilon^{-2} \log N)$ such that the pairwise Euclidean distances between the points are preserved within a factor of $1 \pm \epsilon$, without learning any parameters.

Under our shared-initialization protocol, the pre-trained experts cluster task-specific inputs into highly distinct, well-separated manifolds within the penultimate representation space. The frozen random Gaussian projection acts as a powerful, non-parametric linear dimensionality reduction that preserves these relative geometric clustering distances and semantic manifold separation with high probability, while introducing \textbf{zero new parameters}. Thus, the $d=8$ frozen Gaussian projection acts as a crucial structural regularizer that completely eliminates optimization variance, explaining why it generalizes significantly better than learnable alternatives under tight data limits.

\subsection{Alternative: Unsupervised PCA Projections}

Utilizing the top $d=8$ principal components from the training split (PCA) is an alternative non-parametric approach. While PCA projects along the directions of maximum variance, we find that the classification performance is virtually identical to the random Gaussian projection (e.g., $52.70\% \pm 5.11\%$ on TinyCNN-4 Cross-Domain). Because the penultimate representations of our experts are already highly structured and low-rank due to pre-training, a random projection is sufficient to capture the task-separating directions, rendering the extra computational overhead of PCA redundant. This validates our choice of the random Gaussian projection as a highly elegant, computationally efficient, and statistically robust regularizer.

\section{Conceptual Frontiers and Future Extensions in Dynamic Merging}
\label{app:conceptual_frontiers}

In this section, we expand upon the core theoretical and empirical findings of our paper to outline three promising conceptual frontiers for future dynamic model-merging research. These pathways directly address the core boundaries identified in our deconstruction of the paradigm.

\subsection{Continuous-Time Dynamic Routing (Neural ODEs)}

In deep multi-layer architectures, layer-wise dynamic routing coefficients can exhibit high variance across successive layers, which can exacerbate representation alignment issues and activation drift. To mitigate this optimization noise on scarce calibration budgets, we propose parameterizing the layer-wise routing trajectory as a continuous-time dynamical system.

Specifically, rather than learning independent routing parameters per layer, we can model the coefficient trajectory $\lambda(z)$ across the network depth $z \in [0, 1]$ as the solution to a Neural Ordinary Differential Equation (Neural ODE) \cite{chen2018neural}:
\begin{equation}
    \frac{d\lambda(z)}{dz} = f_{\theta}(\lambda(z), \psi(x))
\end{equation}
where $f_{\theta}$ is a shallow parameter-efficient neural network. Solving this initial value problem from $z=0$ (input layer) to $z=1$ (output layer) generates a smooth, spatially regularized trajectory of merging coefficients. This formulation enforces a continuous topological constraint on the routing space, drastically reducing the parameter footprint and smoothing routing transitions, which can mitigate parameter-variance constraints in deeply stacked physical models.

\subsection{Alternative Non-Linear Parameter Fusion Operators}

We deconstructed the $47\%$ Oracle performance gap on TinyCNN-4 by demonstrating that linear weight-space interpolation acts as a low-pass filter on local spatial convolutional kernels, destroying high-frequency edge and texture responses. To bridge this gap, future research must move beyond the simple linear convex combinations that dominate current literature. We propose exploring two concrete mathematical formulations for non-linear parameter blending operators:

\paragraph{1. Spline-Based Parameter Interpolation.}
Let $W_1^{(l)}, W_2^{(l)}$ be the pre-trained weights of Task 1 and Task 2 experts at layer $l$, and let $\lambda^{(l)}(x) \in [0, 1]$ be our learned BSigmoid routing coefficient. Instead of a linear convex combination $W_{merged}^{(l)}(\lambda) = (1 - \lambda) W_1^{(l)} + \lambda W_2^{(l)}$, we can define weight blending via a learned non-linear cubic Hermite or B\'ezier spline:
\begin{equation}
    W_{merged}^{(l)}(\lambda) = (1 - \lambda)^3 W_1^{(l)} + 3(1-\lambda)^2 \lambda C_1^{(l)} + 3(1-\lambda)\lambda^2 C_2^{(l)} + \lambda^3 W_2^{(l)}
\end{equation}
where $C_1^{(l)}, C_2^{(l)}$ are trainable low-rank control matrices of identical dimensions, initialized to lie on the straight interpolation line:
\begin{equation}
    C_1^{(l)} = \frac{2}{3}W_1^{(l)} + \frac{1}{3}W_2^{(l)}, \quad C_2^{(l)} = \frac{1}{3}W_1^{(l)} + \frac{2}{3}W_2^{(l)}
\end{equation}
During the few-shot calibration phase, these control matrices are optimized to curve and bend the interpolation trajectory. This allows the routing pathway to bypass non-convex, high-loss parameter-space energy barriers, preserving the structural integrity and local spatial frequencies of convolutional filters.

\paragraph{2. Coordinate-Based MLP Weight Generators (Hypernetworks).}
Alternatively, we can express the merged weight as a base static combination plus a dynamic, non-linear spatial delta generated by a lightweight hypernetwork $H_\phi^{(l)}$ conditioned on our low-dimensional routing state:
\begin{equation}
    W_{merged}^{(l)}(\lambda, x) = \bar{W}^{(l)} + H_\phi^{(l)}(\lambda \cdot \psi(x))
\end{equation}
where $\bar{W}^{(l)} = 0.5(W_1^{(l)} + W_2^{(l)})$ is the static uniform average, and $\psi(x) \in \mathbb{R}^d$ is our random Gaussian projected routing state. To remain parameter-efficient and prevent overfitting on our scarce calibration budget, the hypernetwork $H_\phi^{(l)}$ is formulated using low-rank weight-factorization matrices:
\begin{equation}
    H_\phi^{(l)}(\lambda \cdot \psi(x)) = B^{(l)} \cdot \sigma\left(W_{hyper}^{(l)} \cdot (\lambda \cdot \psi(x))\right) \cdot A^{(l)}
\end{equation}
where $A^{(l)} \in \mathbb{R}^{r \times d_{in}}$ and $B^{(l)} \in \mathbb{R}^{d_{out} \times r}$ are frozen random low-rank projections (with rank $r \ll d_{in}, d_{out}$), and $W_{hyper}^{(l)} \in \mathbb{R}^{r \times d}$ is the only trainable parameters. By mapping the dynamic routing state $\lambda \cdot \psi(x)$ directly to non-linear parameter deltas, we preserve high-frequency filter responses without requiring full expert ensembling.


\subsection{PEFT-Level Physical Scale-Up on Vision Transformers}

In Appendix~\ref{app:peft_adapter_collapse_hypothesis} and Appendix~\ref{app:systems_scaling}, we showed that dynamic model merging is highly practical primarily when restricted to low-rank Parameter-Efficient Fine-Tuning (PEFT) modules, such as LoRA adapters, due to systems-level memory-bandwidth limits. We hypothesized that routing over LoRA adapters would exhibit an even deeper spatial specialization (lower SVD Collinearity Ratios) than full-parameter backbones.

A crucial and direct empirical extension of our work is the physical verification of this PEFT adapter collapse hypothesis on large-scale Vision Transformers (e.g., ViT-B/16 CLIP). Evaluating this setup on standard, high-capacity vision benchmarks---such as Stanford Cars \cite{krause20133d}, Oxford Flowers \cite{nilsback2008automated}, and CUB-200 \cite{wah2011caltech}---would allow us to:
\begin{enumerate}
    \item Measure the SVD Collinearity Ratio of learned routing coefficients over specialized LoRA adapters.
    \item Confirm if PEFT-level dynamic merging avoids the massive parameter-space interference and clashing of full weight matrices.
    \item Verify our memory-bandwidth and batch latency scaling projections on physical multi-GPU hardware.
\end{enumerate}
This physical scale-up would bridge the gap between our controlled, calibrated diagnostic sandbox and large-scale, real-world multi-task adaptation.
